\section{Research Questions}\label{sec:research-questions}

The central question guiding this dissertation is whether a small, offline,
unattended language model can be made reliable enough to drive an automated
security copilot once its output is subject to deterministic validation.
From this central question three concrete research questions are derived:

\begin{enumerate}
\item \textbf{RQ1 (pipeline value).} How much does a stacked validation
pipeline improve structured intent-recognition accuracy and flaw catching over
raw language-model output?
\item \textbf{RQ2 (layer contribution).} Which individual guard layer
contributes the most to the flaw-catch rate, and which layers interact
destructively when combined?
\item \textbf{RQ3 (generalisability).} Do the guardrails generalise across two
different offline language-model backends, or is the measured safety an
artefact of a single model?
\end{enumerate}

The study is conducted as an empirical, measurement-based investigation rather
than a hypothesis-testing one: the outcomes of RQ1--RQ3 are quantified from the
committed experiment record reported in Chapter~\ref{ch:results}, and no
statistical significance claim is attached to the single-run measurements
(see Section~\ref{sec:scope-limitations-assumptions}).